\documentclass[11pt,a4paper]{article}
\usepackage[margin=25mm]{geometry}
\usepackage{amsmath,amssymb,amsthm,booktabs}
\newtheorem{lemma}{Lemma}\newtheorem{theorem}{Theorem}
\title{Three maximal overgroups and the order-28 column}
\author{Human Compression project}\date{}
\begin{document}\maketitle
Let $G=O'N$. We recover the entire order-28 column of the fixed
Connor--Leemans Table~1 by counting intersections of three maximal
overgroups. The complete subgroup roster and previously acquired
character evaluations are explicit borrowed premises. Row numbers
below refer to that roster. Its cyclic subgroup of order 28 is row 466;
following the complete immediate-overgroup lists shows that its only
maximal overgroup classes are rows 2, 3 and 7. The first two are copies of
$L_3(7):2$, with the graph automorphism as the outer involution. The
third consists of the centralizers of involutions.

\begin{lemma}\label{normalizer}
Put $M=L_3(7):2$. If $x\in M$ has order 28 and $U=\langle x^4\rangle$,
then $U$ is generated by a transvection and $|N_M(U)|=8232$.
If $M\leq G$, then $N_M(U)=N_G(U)$.
\end{lemma}
\begin{proof}
The element $x^4$ lies in $L_3(7)$ and has order 7. An element of order 7
has a unipotent lift to $\mathrm{SL}_3(7)$: any scalar seventh power can
be removed by the unique appropriate scalar in its center of order 3.
There are two nonidentity Jordan types, a transvection and a regular
unipotent element. For a regular unipotent matrix $1+N$, every commuting
matrix is $aI+bN+cN^2$. The determinant-one condition is $a^3=1$;
after division by the scalar center the centralizer has order 49.
A projective commutator scalar must be1, since conjugate unipotent
matrices have the same eigenvalue 1. Its centralizer in $M$ consequently has order 49 or 98. It cannot
contain an element of order 28. Since $x$ centralizes $x^4$, the regular
case is excluded.

Choose coordinates with $U=\langle I+E_{13}\rangle$. A matrix
normalizing $U$ preserves both the image line of $E_{13}$ and its
kernel plane. Conversely, every invertible upper triangular matrix
conjugates $E_{13}$ to a nonzero scalar multiple. Thus the normalizer
in $L_3(7)$ is the projective image of the determinant-one upper
triangular group. Its order is
\[
                         7^3\,6^2/3=4116.
\]
Inverse transpose followed by conjugation with an antidiagonal Weyl
matrix also normalizes $U$, and lies outside $L_3(7)$. This doubles
the normalizer to 8232. The roster has exactly two cyclic order 7
classes, with normalizer orders 294 and 8232. Since $N_M(U)$ is contained
in $N_G(U)$, the first is impossible and the second gives equality.
\end{proof}

The roster identifies this global normalizer as row 22. Every row 22
subgroup belongs to precisely one maximal subgroup in each of rows 2
and 3: the direct incidence calculation is
\[
                       122760\cdot456=55978560,
\]
where 55978560 is the class length of row 22. This is a count of actual
incidences, not a product along paths in the subgroup diagram.

\begin{lemma}\label{three}
Every element $x$ of order 28 in $G$ lies in exactly three maximal
subgroups $A,B,C$, in rows 2, 3 and 7 respectively. Their intersections
have rows and orders
\[
\begin{array}{c|rrrr}
 &A\cap B&A\cap C&B\cap C&A\cap B\cap C\\\hline
\text{row}&22&65&66&238\\
\text{order}&8232&1344&1344&168.
\end{array}
\]
\end{lemma}
\begin{proof}
The cyclic 28 ancestor list supplies a projective maximal overgroup
for every such $x$. First consider one containing $x$.
Lemma~\ref{normalizer} forces it to contain the fixed subgroup
$N_G(\langle x^4\rangle)$ of row 22. The preceding incidence count
therefore gives at most one subgroup of each projective maximal class.
Conversely, both maximal overgroups of that row 22 subgroup contain
$x$, since $x$ normalizes $\langle x^4\rangle$.

The accepted centralizer criterion says that if the cyclic 2-part of
$x$ is a Sylow 2-subgroup of $C_G(x)$, then $C_G(x)$ has a unique
involution. The primary class data give $|C_G(x)|=28$ for both
order 28 classes. Consequently the only involution centralizer
containing $x$ is $C=C_G(t)$, where $t=x^{14}$. The complete
cyclic 28 overgroup list excludes every other maximal class.

The subgroup $N_G(\langle x^4\rangle)$ is maximal in $A$ and is
contained in the distinct maximal subgroup $B$. Hence it is $A\cap B$.
The involution $t$ lies in the simple subgroup of $A$, since its image
in the quotient of order 2 is the fourteenth power of the image of $x$.
The accepted inner-involution centralizer calculation in $L_3(7):2$
gives $|C_A(t)|=1344$. Its roster class is 65; for $B$ it is 66.
Thus $A\cap C=C_A(t)$ and $B\cap C=C_B(t)$.

It remains to compute
\[
 A\cap B\cap C=N_{C_A(t)}(\langle x^4\rangle).
\]
The group $C_A(t)$ is the block group
$(\mathrm{GL}_2(7)/C_3):2$ used in the accepted centralizer proof.
Every order 7 subgroup in it is conjugate to the upper unipotent
subgroup of its $2$ by $2$ block. Indeed an order 7 element lifts,
after removal of a central scalar, to a nonidentity unipotent matrix,
and all such subgroups in $\mathrm{GL}_2(7)$ are conjugate.
Its normalizer in $\mathrm{GL}_2(7)$ is the upper triangular group
of order$7\cdot6^2=252$. Dividing by the scalar subgroup of order 3
and adjoining inverse transpose followed by the $2$ by $2$ Weyl
matrix gives order 168. The intersection contains $x$, so the complete
cyclic 28 ancestor list identifies its class as row 238.
\end{proof}

For clarity, the block normalizer calculation can also be checked
without naming its structure. In the retained 1344-element permutation
group, choose a displayed element$u$ of order 7 and test
$g^{-1}ug\in\langle u\rangle$ for every element$g$. The returned set
has 168 elements, and its displayed generators span precisely that set.
This literal calculation is included in the certificate.

\begin{theorem}\label{formula}
Let $x_L(p,q)$ count ordered pairs $(r,s)$ in $L$ with orders $p,q$
and with $rs$ an involution, and let $e_G(p,q)$ count those generating $G$.
For every$p$ in the source's order list,
\begin{align*}
 e_G(p,28)={}&x_G(p,28)-\sum_{M}[G:M]x_M(p,28)\\
 &+55978560x_{T}(p,28)
  +2\cdot342868680x_{H}(p,28)
  -2742949440x_{K}(p,28),
\end{align*}
where the sum runs over the thirteen maximal conjugacy classes,
and $T,H,K$ are the explicit groups of orders 8232, 1344, 168 above.
\end{theorem}
\begin{proof}
Every nongenerating pair lies in a maximal subgroup. For a pair whose
second member has order 28, Lemma~\ref{three} gives at most the three
specified maximal subgroups. Inclusion--exclusion therefore stops at
their triple intersection.

We check that this counts subgroup classes, rather than merely
possible intersection types. Every row 22 subgroup has unique actual
overgroups in rows 2 and 3, and their intersection is that row 22 subgroup.
Every row 65 subgroup has unique actual overgroups in rows 2 and 7;
row 66 has unique actual overgroups in rows 3 and 7. The direct incidence
identities for these statements are
\[
 122760\cdot2793=2857239\cdot120=342868680.
\]
All these subgroups contain elements of order 28. A row 238 subgroup
also contains such an element. Its complete immediate-overgroup list
and direct counts give a unique row 22 overgroup and unique overgroups
in rows 65 and 66:
\[
 55978560\cdot49=342868680\cdot8=2742949440.
\]
Using an order 28 element in the row 238 subgroup aligns these overgroups
with the same three maximal subgroups of Lemma~\ref{three}. Their
intersection has order 168 and hence is exactly the given subgroup.
Thus each intersection family is counted once with the class length
in the displayed formula. Rows 65 and 66 are isomorphic block groups,
so their two identical intrinsic contributions are added. This proves
the formula, including possible pairs with equal generator orders.
\end{proof}

The finite certificate uses the retained 114-point projective action.
It constructs the upper triangular group and its graph extension,
reuses the block group, and extracts its order 7 normalizer. Literal
closure and cardinality checks identify the three groups. In each,
the elements of order 28 form two conjugacy classes, of total sizes
588, 96 and 12 respectively. For each class representative, the reader
tests every partner and multiplies by the class size. The nonzero
intrinsic counts are
\[\begin{array}{c|rrr}
p&x_T(p,28)&x_H(p,28)&x_K(p,28)\\\hline
2&16464&2688&336\\
3&0&672&0\\
4&4704&0&0\\
6&0&3360&0\\
7&0&672&0\\
8&0&1344&0\\
14&98784&672&0\\
16&0&5376&0\\
28&24108&96&12\\
\end{array}\]

These computations use no ambient enumeration of $O'N$ and no group-order
or character-table software call. The independent reader uses the
opposite permutation-composition convention; the order of a product
is unchanged when its two factors are reversed.

Substitution into Theorem~\ref{formula}, and division by
$|\operatorname{Aut}(G)|=921631011840$, gives
\[\begin{array}{c|rr}
(p,q)&e_G(p,q)/|\operatorname{Aut}(G)|&\text{previously unpaid}\\\hline
(3,28)&20&\text{yes}\\
(4,28)&370&\text{yes}\\
(5,28)&560&\text{no}\\
(6,28)&1370&\text{yes}\\
(7,28)&2214&\text{yes}\\
(8,28)&6276&\text{no}\\
(10,28)&4992&\text{no}\\
(11,28)&9202&\text{no}\\
(12,28)&2796&\text{yes}\\
(14,28)&3616&\text{yes}\\
(15,28)&4496&\text{no}\\
(16,28)&25316&\text{no}\\
(19,28)&15978&\text{no}\\
(20,28)&10292&\text{no}\\
(28,28)&7246&\text{yes}\\
(28,31)&6658&\text{no}\\
\end{array}\]

For $p=31$ the last row exchanges the generator roles; this is an
automorphism-equivariant bijection of the two ordered-pair sets, not a
quotient by reversal. All sixteen values agree with the fixed source. Seven were previously
unpaid. There is no division by duality: the source counts automorphism
orbits of ordered pairs, and the automorphism action on generating
pairs is free. The character evaluations, complete subgroup roster,
projective maximal identifications and involution centralizer data are
borrowed inputs. The contribution is the canonical three-overgroup
argument, its exact inclusion--exclusion formula, and the literal
small-group certificate. No new chirality claim or whole-source
completion follows from this column alone.

\begin{thebibliography}{9}
\bibitem{CL} T. Connor and D. Leemans, \emph{Algorithmic enumeration
of regular maps}, Ars Math. Contemp. 10 (2016), 211--222, Table 1.
\bibitem{roster} T. Connor and D. Leemans, \emph{The subgroup lattice
of $O'N$}, fixed 23-page roster. The exact native pages and class
identifications used here are bound with the certificate.
\bibitem{prior} Human Compression project, \emph{Central involutions
and five exact map counts} and \emph{Four O'Nan map counts from one small group}, accepted prior proof components; exact editions and
distinct reviews are imported.
\end{thebibliography}
\end{document}
